% !TEX TS-program = pdflatex
% CV ciblé — BHCT : Développeur Full Stack Java / ReactJS
\documentclass[a4paper,10pt]{article}
\usepackage[utf8]{inputenc}
\usepackage[T1]{fontenc}
\usepackage[scaled=0.94]{helvet}
\renewcommand{\familydefault}{\sfdefault}
\usepackage[left=1.0cm,right=1.0cm,top=0.5cm,bottom=0.4cm]{geometry}
\usepackage{enumitem}
\usepackage{titlesec}
\usepackage{xcolor}
\usepackage[hidelinks]{hyperref}
\definecolor{navy}{RGB}{23,54,93}
\pagestyle{empty}
\setlength{\parindent}{0pt}
\setlength{\parskip}{0pt}
\linespread{0.72}
\hypersetup{pdftitle={CV - Baha Eddine Belhaj Mouldi - Développeur Full Stack Java / ReactJS},pdfauthor={Baha Eddine Belhaj Mouldi},pdfsubject={Java, Spring Boot, ReactJS, APIs REST, Microservices}}
\titleformat{\section}{\normalsize\bfseries\color{navy}}{}{0pt}{\MakeUppercase}[{\vspace{1pt}\color{navy}\hrule height 0.6pt}]
\titlespacing*{\section}{0pt}{3pt}{1.5pt}
\setlist[itemize]{leftmargin=1.1em,label=\textbullet,itemsep=0.1pt,topsep=0.2pt,parsep=0pt}
\newcommand{\cventry}[4]{\textbf{#1} \hfill \textbf{#4}\\[0.5pt]#2{\ifx\relax#3\relax\else, #3\fi}\par\vspace{1pt}}
\begin{document}
\begin{center}
{\LARGE\bfseries\color{navy} Baha Eddine Belhaj Mouldi}\\[3pt]
{\large Développeur Full Stack Java / ReactJS --- Spring Boot \textbar React \textbar Microservices}\\[4pt]
Berges du Lac 2, Tunis \quad|\quad +216 55 128 982 \quad|\quad \href{mailto:bahaeddine.belhajmouldi@gmail.com}{bahaeddine.belhajmouldi@gmail.com}\\[1pt]
\href{https://www.linkedin.com/in/bahamouldi}{linkedin.com/in/bahamouldi} \quad|\quad \href{https://github.com/bahamouldi}{github.com/bahamouldi} \quad|\quad \href{https://bahamouldi.github.io/portfolio/}{bahamouldi.github.io/portfolio}
\end{center}
\vspace{2pt}
\section{Profil}
Ingénieur diplômé en génie informatique (ENIT, 65e/600+) avec 3+ ans d'expérience en développement Full Stack Java/Spring Boot et ReactJS. Conception et réalisation de bout en bout d'applications web robustes : analyse des besoins, architecture technique, développement backend Java (Spring Boot, Spring Security, Spring Data), frontend React/TypeScript, tests unitaires, revues de code et mise en production. Expérience concrète en microservices, conteneurisation Docker et pipelines CI/CD (Jenkins). Méthodologies Agile/Scrum. Excellent niveau de français, oral et écrit.

\section{Compétences clés}
\textbf{Back-end \& Java} : Java 8/11/17, Spring Boot 3.x, Spring Security (JWT, OAuth2), Spring Data JPA/Hibernate, REST APIs, microservices, validation de beans, gestion des erreurs\\[0.5pt]
\textbf{Front-end \& React} : ReactJS, React Router, hooks (useState, useEffect, useMemo), composants fonctionnels, TypeScript, HTML5/CSS3, Tailwind CSS, interfaces responsive\\[0.5pt]
\textbf{Base de données} : SQL avancé, PostgreSQL, MySQL, Elasticsearch (indexation, requêtes), modélisation de données, JDBC, JPA\\[0.5pt]
\textbf{Outils \& Technologies} : Git, Maven, Jenkins (pipelines CI/CD), Docker/Docker Compose, Elasticsearch, Jira, Confluence\\[0.5pt]
\textbf{Qualité \& Tests} : JUnit 5, Mockito, tests unitaires et d'intégration, Swagger/OpenAPI, revue de code, architecture modulaire\\[0.5pt]
\textbf{Méthodologies} : Agile/Scrum, sprints, rétrospectives, Kanban, rédaction de documentation technique

\section{Expérience professionnelle}
\cventry{Développeur Full Stack Java / ReactJS --- Freelance}{Digital Power Consulting}{Tunisie}{01/2024 -- 05/2026}
\begin{itemize}
  \item Réalisation de bout en bout de 5+ applications web full stack : backend Java/Spring Boot (Spring Security, Spring Data, REST APIs), frontend ReactJS/TypeScript, bases de données PostgreSQL/MySQL.
  \item Architecture microservices : découplage de services, communication REST, scalabilité horizontale, circuits breakers.
  \item Sécurité applicative : authentification JWT/OAuth2, gestion des rôles/permissions, codage sécurisé conforme OWASP.
  \item Conteneurisation Docker, CI/CD Jenkins, déploiement Kubernetes. Tests unitaires JUnit 5/Mockito, revues de code.
\end{itemize}
\cventry{Développeur Full Stack --- Stagiaire PFE}{Digital Power Consulting}{Tunisie}{01/2026 -- 05/2026}
\begin{itemize}
  \item Application full-stack : backend Java/Spring Boot (JPA/Hibernate, REST APIs, Spring Security), frontend ReactJS/TypeScript, PostgreSQL, Docker.
  \item Dashboard interactif temps réel : visualisation de données, graphiques dynamiques, interface responsive mobile-first.
  \item Architecture technique documentée, choix d'architecture validés avec l'équipe, mise en production via pipeline CI/CD Jenkins.
\end{itemize}
\cventry{Stagiaire Cybersécurité --- Détection menaces SAP}{Streamlink}{Tunisie}{07/2025 -- 08/2025}
\begin{itemize}
  \item Scripts Python (PyRFC), dashboards ELK (Elasticsearch, Kibana), scoring de risque automatisé pour 7 transactions critiques SAP.
\end{itemize}
\cventry{Stagiaire Sécurité Réseau}{Tunisie Telecom}{Tunisie}{07/2024}
\begin{itemize}
  \item Dashboards Elastic/Kibana, analyse trafic réseau, détection d'anomalies, identification de 10+ vulnérabilités.
\end{itemize}

\section{Projets clés}
\cventry{Architecture Microservices --- Kafka + Spring Boot}{Java, Spring Boot, ReactJS, PostgreSQL, Docker, Kubernetes, Kafka}{}{2026}
\begin{itemize}
  \item 4 microservices Java/Spring Boot découplés, API REST inter-services, Kafka (topics ordonnés, exactly-once), frontend ReactJS, PostgreSQL par service, Docker + Kubernetes.
  \item Spring Security (JWT), Spring Data JPA, tests unitaires JUnit 5/Mockito, documentation Swagger/OpenAPI.
\end{itemize}
\cventry{IDTS --- Plateforme de Gestion Sécurisée}{Java, Spring Boot, ReactJS, PostgreSQL, Docker, Kubernetes}{}{2024}
\begin{itemize}
  \item Backend Spring Boot (JPA/Hibernate, REST APIs, validation, gestion des erreurs), frontend ReactJS (CRUD, hooks, routes protégées), PostgreSQL, Docker, Kubernetes.
\end{itemize}

\cventry{Pipeline Streaming Temps Réel --- Kafka + PySpark}{PySpark, Kafka, PostgreSQL, Docker Compose, Flask}{}{2024}
\begin{itemize}
  \item Architecture événementielle complète (6 services Docker Compose) : Kafka Producer → consumer PySpark Structured Streaming → fenêtrage 30 min + watermark → PostgreSQL via JDBC.
\end{itemize}
\cventry{Dashboard Assurance Revenus --- Tunisie Telecom}{Python, Streamlit, PostgreSQL, Plotly, RAG}{}{2026}
\begin{itemize}
  \item Pipeline ETL 900K+ enregistrements : data warehouse schéma étoile, dashboards interactifs, chatbot RAG, modèles de prévision.
\end{itemize}

\section{Formation}
\cventry{Diplôme d'Ingénieur en Génie Informatique}{École Nationale d'Ingénieurs de Tunis (ENIT)}{Tunisie}{09/2023 -- 06/2026}
\begin{itemize}
  \item Classé 65e/600+ au concours national d'entrée. PFE : Architecture microservices Java/Spring Boot avec Kafka (86/100).
\end{itemize}
\cventry{Classes préparatoires (Physique-Technologie)}{Institut Préparatoire aux Études d'Ingénieurs, El Manar}{Tunisie}{09/2021 -- 06/2023}

\section{Certifications}
Git \& GitHub --- 365 Data Science (2024) ; Réseaux --- Cisco (2025) ; Fondements du Deep Learning --- NVIDIA (2025) ; Machine Learning --- NVIDIA (2024).

\section{Langues}
Français (Courant, écrit et oral) \quad|\quad Anglais (B2) \quad|\quad Arabe (Natif)

\end{document}
